\section{Practical Analysis}
\label{sec:practical}

\subsection{R:r = 30:1 at $d = 308$}

The IOT implementation uses $R/r = 30$ with $d = 154$ torus pairs (308 scalar dimensions). We evaluate the critical scaling:
\begin{align}
    \left(\frac{r}{R}\right)_{\text{actual}} &= \frac{1}{30} = 0.0\overline{3} \\[4pt]
    \left(\frac{r}{R}\right)_{\text{critical}} &= \frac{\ln 154}{2 \cdot 154} = \frac{5.037}{308} = 0.01636
\end{align}

The actual ratio exceeds the critical ratio by a factor of $0.0333 / 0.0164 = 2.04$, placing the system firmly in the \textbf{expansion regime}.

Computing the contrast:
\begin{align}
    \left(\frac{31}{29}\right)^{77} &= e^{77 \ln(31/29)} = e^{77 \times 0.0667} = e^{5.14} \approx 170 \\[4pt]
    C(154) &\approx \frac{170}{\sqrt{154}} = \frac{170}{12.41} \approx 13.7
\end{align}

\begin{remark}
A contrast ratio of 13.7:1 means that for uniformly distributed points on the 308-dimensional IOT, the farthest neighbor is on average $13.7\times$ farther than the nearest neighbor. For comparison, Euclidean $L_2$ at 308 dimensions yields contrast approaching 0, meaning nearest and farthest neighbors are indistinguishable.
\end{remark}

\subsection{Alternative $r/R$ Scalings}

Table~\ref{tab:scalings} evaluates several $r/R$ scalings at $d = 154$ torus pairs.

\begin{table}[h]
\centering
\footnotesize
\begin{tabular}{@{}llrl@{}}
\toprule
Scaling & $R/r$ at $d\!=\!154$ & $C(154)$ & Regime \\
\midrule
$r/R = 1/d$ & 154 & $\approx 0.08$ & Contraction \\
$r/R = \ln d/(2d)$ & 61.2 & $\approx 1.0$ & Critical \\
$r/R = 1/30$ (fixed) & 30 & $\approx 13.7$ & Expansion \\
$r/R = 1/\sqrt{d}$ & 12.4 & $\approx 1.6 \times 10^4$ & Strong expansion \\
$r/R = 1/\ln d$ & 5.04 & $\approx 10^{14}$ & Degenerate \\
\bottomrule
\end{tabular}
\caption{Contrast at $d = 154$ torus pairs for various scalings.}
\label{tab:scalings}
\end{table}

The $r/R = 1/\sqrt{d}$ scaling yields $C(d) \sim e^{\sqrt{d}}/\sqrt{d}$, which grows superexponentially. This is theoretically interesting but practically problematic: when contrast is too large, the IOT distance is dominated by the binary path choice rather than angular proximity. Fine-grained similarity is lost.

\subsection{Optimal Operating Point}

The ideal operating point balances two requirements:
\begin{enumerate}
    \item \textbf{Expansion}: $r/R$ must exceed the critical threshold $\ln(d)/(2d)$ to ensure the curse is reversed.
    \item \textbf{Fine-grained discrimination}: $r/R$ must not be so large that the binary path choice overwhelms the continuous angular distance. The distance metric should reflect geometric proximity, not merely which dimensions chose the involuted path.
\end{enumerate}

We propose the operating range:
\begin{equation}
    \frac{r}{R} = c \cdot \frac{\ln d}{d}, \quad c \in (1, 3)
\end{equation}

At $c = 1$, the system is exactly at criticality. At $c = 3$, the contrast is $C(d) \sim d^2/\sqrt{d} = d^{3/2}$, which grows polynomially---substantial discrimination without degenerate dominance of the path-choice mechanism.

For $d = 154$ torus pairs, $c = 2$ gives $R/r \approx 30.6$, remarkably close to the $R{:}r = 30{:}1$ ratio used in the IOT implementation.

\subsection{Dimension-Adaptive Formulation}

For applications requiring a single formulation that works across varying dimensionalities, we propose:
\begin{equation}
    R(d) = \frac{d}{c \cdot \ln d} \cdot r
    \label{eq:adaptive}
\end{equation}
with $c = 2$ as the default operating point. This ensures the system operates at $2\times$ criticality regardless of $d$, maintaining meaningful contrast without degenerate behavior.

\begin{table}[h]
\centering
\footnotesize
\begin{tabular}{@{}rcrc@{}}
\toprule
$d$ (pairs) & Scalar dims & $R/r$ at $c\!=\!2$ & $C(d)$ \\
\midrule
16 & 32 & 2.9 & 4.0 \\
64 & 128 & 7.7 & 8.0 \\
154 & 308 & 30.6 & 13.7 \\
512 & 1024 & 41.1 & 22.6 \\
2048 & 4096 & 134.2 & 45.3 \\
\bottomrule
\end{tabular}
\caption{Dimension-adaptive $R/r$ and contrast at $c = 2$.}
\label{tab:adaptive}
\end{table}

The contrast grows as $\Theta(\sqrt{d})$---polynomial, not exponential---providing steadily increasing discrimination without divergence.
